\documentclass[12pt]{article}
\usepackage[utf8]{inputenc}
\usepackage[spanish]{babel}

\usepackage[backend=biber, style=numeric]{biblatex}
\addbibresource{ref.bib}	

\usepackage[margin=1in]{geometry}		% Margen
\usepackage{booktabs}					% Tablas
\usepackage{graphicx}       			% Insertar imgs
\usepackage[style=ddmmyyyy]{datetime2}       % Fijar formato día
\usepackage{amsthm, amsfonts, amsmath, dsfont, amssymb}         % For Math
\usepackage{fancyhdr}       % For fancy header/ footer
\usepackage{cancel}         % Cancelar expresiones
\usepackage[ruled,vlined]{algorithm2e} % To write algorithms
\usepackage{tikz}          % Visuales
\usetikzlibrary{shapes.geometric} % Más figuras

\usepackage{xcolor} % Personalizar color

\usepackage{eso-pic} % Fondo portada





%%%%%%%%%%%%%%%%%%%%%%
% info basica
\newcommand{\myname}{- Carlos Gabriel Barrera Casasola,\\ - Diego Marquez Franco,\\ - Pablo Rogelio Ramírez Alferes}
\newcommand{\myclass}{Un caso degenerado del modelo CIR/Feller}
\newcommand{\mytitle}{Proyecto: Ecuaciones Diferenciales Estocásticas}
% \newcomand{\mydate}{<Date>}

\title{\mytitle}
\author{\myname}
% \date{\mydate}
%%%%%%%%%%%%%%%%%%%%%%



%%%%%%%%%%%%%%%%%%%%%%
% Set up fancy header/footer
\pagestyle{fancy}
\fancyhead[LO,L]{ \textbf{\large{\myclass}} \\  \textbf{\textit{\mytitle}}}
\fancyhead[CO,C]{ \\ \hspace{6.0cm} }
\fancyhead[RO,R]{\textbf{\textit{}} \\ \textit{\DTMtoday}} % \DTMtoday for today, \mydate for custom date
\fancyfoot[LO,L]{}
\fancyfoot[CO,C]{\thepage}
\fancyfoot[RO,R]{}
\renewcommand{\headrulewidth}{0.4pt}
\renewcommand{\footrulewidth}{0.4pt}
\setlength{\headsep}{1.2cm}

%%%%%%%%%%%%%%%%%%%%%%%%%%%%%%%%%%%%%%%%%%%%%%%%%%%%%%%%%%%%%%%%%%%%%%%%%%%
%%%%%%%%%%%%%%%%%%%%%%% Define all premade object %%%%%%%%%%%%%%%%%%%%%%%%%
%%%%%%%%%%%%%%%%%%%%%%%%%%%%%%%%%%%%%%%%%%%%%%%%%%%%%%%%%%%%%%%%%%%%%%%%%%%


%%%%%%% Triangulo para enumerar
\newcommand*\mytriangle[1]{\hspace*{-0pt}\raisebox{0.8ex}{\tikz[baseline=(char.base)]{
			\node[shape=regular polygon, regular polygon sides=5, fill=green!65!gray!40!white, inner sep=0pt, rotate=-90, scale=1.15, anchor=base] (char) {#1};}}
		}



%%%%%%%%%%%%%%%%%%%%%%%%%%%%%%%%%%%%%%%%%%%%%%%%%%%%%%%%%%%%%%%%%%%%%%%%%%%


%%%%%%%%%%%%%%%%%%%%%%%%%% Portada  %%%%%%%%%%%%%%%%%%%%%%%%%%%%%%%%
\newcommand{\marcaDeAgua}{%
	\begin{tikzpicture}[remember picture, overlay]
		\node at (current page.center) {%
			\begin{tikzpicture}
				% Espiral mu
				\foreach \t in {0,12,...,1080} {
					\pgfmathsetmacro{\r}{0.022*\t}
					\pgfmathsetmacro{\x}{\r*cos(\t)}
					\pgfmathsetmacro{\y}{\r*sin(\t)}
					\pgfmathsetmacro{\s}{0.25+0.011*\t}
					\node[scale=\s, gray!50!green!25!white!55, rotate=\t] at (\x,\y) {$\mu$};
				}
				% sigma
				\foreach \t in {0,12,...,1080} {
					\pgfmathsetmacro{\r}{0.02*\t}
					\pgfmathsetmacro{\ang}{\t+180} % Para que no quede al revez
					\pgfmathsetmacro{\x}{\r*cos(\ang)}
					\pgfmathsetmacro{\y}{\r*sin(\ang)}
					\pgfmathsetmacro{\s}{0.25+0.011*\t}
					\node[scale=\s, gray!50!green!25!white!55, rotate=\ang] at (\x,\y) {$\sigma$};
				}
			\end{tikzpicture}%
		};
	\end{tikzpicture}%
}
\AddToShipoutPictureBG*{\marcaDeAgua}


% Para referenciar
\usepackage{hyperref}
\hypersetup{
	colorlinks=true,
	linkcolor=blue,
	filecolor=magenta,      
	urlcolor=cyan,
	citecolor=blue
}


%%%%%%%%%%%%%%%%%%%%%%%%%%%%%%%%%%%%%%%%%%%%%%%%%%%%%%%%%%%%%%%%%%%%%%%%%%%%
\begin{document}
	%%%%%%%%%%%%%%%%%%%%%%%%%%%%%%%%%%%%%%%%%%%%%%%%%%%%%%%%%%%%%%%%%%%%%%%%%%%%%% Portada
	\begin{titlepage}
		\centering
		\vspace*{3cm}
		
		{\Huge \bfseries \myclass \par}
		\vspace{2cm}
		
		{\Large \mytitle \par}
		\vspace{4cm}
		
		{\large \myname \par}
		\vspace{0.5cm}
		
		{\large \DTMtoday \par}
		
		\vfill
		
	\end{titlepage}
	
%%%%%%%%%%%%%% Trabajo


\section*{La ecuación}
Supongamos que $X_t$ es una concentración de elementos de algún fenómeno que decrece con el tiempo. Ejemplo: Decaimiento radiactivo o disolución de un fármaco..

En un intervalo de tiempo $\Delta t$ una fracción $\lambda \Delta t$ $X_t$ desaparece. 
Entonces $$\Delta \approx -\lambda X_t \Delta t$$

Agregando un factor de ruido, digamos que un elemento de lo que representa $X_t$ tiene una probabilidad $\lambda \Delta t$ de desaparecer. Así, $\mathbb{E}[\Delta X_t] = -\lambda X_t \Delta t $ y $\operatorname{Var}(\Delta X_t) = \lambda X_t \Delta t.$
Entonces  $$\sqrt{\operatorname{Var}(\Delta X_t)} =  \sqrt{\lambda X_t \Delta t} = \sqrt{\lambda X_t}$$

Suponiendo cierta normalidad podemos ajustar el termino del ruido al incremento de un browniano $\Delta W_t$ con varianza adecuada, así
$$\Delta X_t = - \lambda X_t \Delta t + \sqrt{ \lambda X_t} \Delta W_t$$

Tomando el limite $\Delta t \to 0$ tenemos la ecuación con la que trabajaremos.

\begin{equation}
	dX_t = -\lambda X_t dt + \sqrt{\lambda X_t } dW_t, \quad \lambda >0
	\label{eq:mainEq}
\end{equation}


Que es un caso particular del modelo CIR con $a=\lambda\ b= 0\ \sigma=\sqrt{\lambda}$ y de la ecuación de Feller con $[params]$.

\section*{Ecuaciones relacionadas}
Observamos que nuestra ecuación \ref{eq:mainEq} es un caso particular y degenerado del modelo \emph{CIR} y de la ecuación de $Feller$.

\subsection*{Cox-Inderson-Ross (CIR)}
Propuesto en 1985, es un modelo para tasas de interés no negativas. 

Está dado por la EDE
$$dX_t = a(b-X_t)dt + \sigma \sqrt{X_t} dW_t, \quad b\in \mathbb{R},\  a,\sigma>0$$

En comparación, la ecuación \ref{eq:mainEq} es el caso donde $a = \lambda, b = 0, \sigma = \sqrt{\lambda}.$


\subsection*{Ecuación de Feller}
Con una gran variedad de aplicaciones por su deriva lineal y difusión dependiente del estado se ha usado en neurobiología, representación de poblaciones y mercados financieros \cite{MasoliverPerello2012}.

Tiene la expresión
$$dX_t = (-\alpha X_t + \beta) dt + k \sqrt{X_t}dW_t,\quad \beta \in \mathbb{R}, \alpha, k >0.$$


Es equivalente al CIR salvo la caracterización de sus parámetros.

La ecuación \ref{eq:mainEq} es un caso particular con $\beta = 0, \alpha = \lambda, k = \sqrt{\lambda}.$


\section*{Análisis de la ecuación}
Tenemos que $b(t,X_t) = -\lambda X_t, \ \sigma(t,X_t)= \sqrt{\lambda X_t} $

\subsection*{Existencia y unicidad}
Veamos si cumple con la condición de Lipschitz y con crecimiento lineal. El tiempo no da problema así que centrémonos en las espaciales.

\subsubsection*{Condición de Lipschitz}
Sean $x,y \geq 0.$
\begin{align*}
	|b(t,x) - b(t,y)|^2 + |\sigma(t, x) - \sigma(t,y)|^2 &= |-\lambda x + \lambda y|^2 + |\sqrt{\lambda x} - \sqrt{\lambda y}|^2 \\
	& = \lambda^2|x-y|^2 + \lambda|\sqrt{x} - \sqrt{y}|^2
\end{align*}

El primer término, correspondiente a $b(t,X_t)$ ya está acotado de la forma deseada, $b$ es Lipschitz. Sin embargo, ¿qué hay de $\sigma(t, X_t)$?\\

Tenemos que $$\frac{d\sigma}{dt} = \frac{\sqrt{\lambda}}{2\sqrt{x}} \xrightarrow[x\to 0]{} \infty$$

Por lo que en 0, los incrementos $|\sigma(t,x)-\sigma(t,y)|$ no pueden ser acotados por diferencias lineales $|x-y|.$

Sin embargo, para $\varepsilon>0$ y $x,y\leq 0$
\begin{align*}
	|\sigma(t,x) - \sigma(t, y)|^2 &= \lambda|\sqrt{x}-\sqrt{y}|^2 \\
	&= \lambda (\frac{|\sqrt{x}-\sqrt{y}|(\sqrt{x}+\sqrt{y})}{\sqrt{x}+\sqrt{y}})^2 \\
	&= \lambda (\frac{|x-y|}{\sqrt{x} + \sqrt{y}})^2
\end{align*}
Ahora, $0<\varepsilon \leq x,y\ \implies 0< 2\sqrt{\varepsilon} \leq \sqrt{x} + \sqrt{y}$

Tomando los inversos se sigue que
\begin{align*}
	|\sigma(t,x) - \sigma(t, y)|^2 & \leq  \lambda (\frac{|x-y|}{2\sqrt{\varepsilon}})^2\\
	&= \frac{\lambda}{4\varepsilon} |x-y|^2
\end{align*}

Habiendo acotad ambos términos de la manera deseada, nuestra ecuación cumple la condición de Lipschitz en $[\varepsilon, \infty)\ \forall \varepsilon > 0.$


\subsubsection*{Condición de crecimiento lineal}
Sea $x\geq 0$.

\begin{align*}
	|b(t,x)|^2 + |\sigma(t,x)|^2 &= |-\lambda x|^2 +|\sqrt{- \lambda x}|^2\\
	&= \lambda^2 |x|^2 + \lambda x
\end{align*}
¿Existirá algún $K$ tal que $\lambda^2 x^2 + \lambda x \leq K(1+x^2)$?\quad  (Recordamos que $X_t \geq 0$)

Esto es equivalente a si $h(x) = \frac{\lambda^2 x^2 + \lambda x}{1+x^2} \leq K.$\\

Desarrollamos buscando un máximo de $h$ positivo.
\begin{align*}
	h(x) & = \frac{\lambda^2 x^2 + \lambda x}{1+x^2} + \frac{\lambda^2}{1+x^2} - \frac{\lambda^2}{1+x^2}\\
	 &= \frac{\lambda^2 (1+x^2) + \lambda x}{1+x^2} - \frac{\lambda^2}{1-x^2}\\
	 &= \lambda^2 + \frac{\lambda x}{1+x^2} - \frac{\lambda^2}{1+x^2}
\end{align*}

Así derivando e igualando a 0
\begin{align*}
	0 = h'(x)& = \lambda \frac{1+x^2 -2x^2}{(1+x^2)^2} - \lambda^2 \frac{-2x}{(1+x^2)^2}\\
	&= \frac{\lambda (1-x^2)}{(1+x^2)^2} + \frac{\lambda^2 2x}{(1+x^2)^2}
\end{align*}

Que ocurre si y solo si
\begin{align*}
	0 & =-[ (1-x^2) +2 \lambda x]\\
	&=  x^2 -2\lambda x - 1
\end{align*}

Resolviendo por la formula general de segundo grado
$$x^* = \lambda \pm \sqrt{\lambda^2 +1}  $$

Pero $\sqrt{\lambda^2 +1 } > \sqrt{\lambda^2} = |\lambda| = \lambda$

Por lo que tomamos $x^*_+ = \lambda + \sqrt{\lambda^2 +1}$ para que $x^*_+$ esté donde la ecuación diferencial estocástica está definida.

Ahora veamos que en $x^*_+$ hay un máximo.

Nos fijamos en los signos de la derivada.
\begin{align*}
	h'(x) &= \lambda \frac{1+x^2 -2x^2}{(1+x^2)^2} - \lambda^2 \frac{-2x}{(1+x^2)^2} \\
	& =\lambda \frac{1-x^2 + 2 \lambda x}{(1+x^2)^2}
\end{align*}

El signo del denominador siempre es positivo.

Para el numerador notamos que se trata de una parábola, como $\lambda > 0$, abre hacia abajo pues el signo del término de segundo grado es negativo. Y además, el signo es positivo en $(x^*_-, x^*_+)$

De esta forma $sgn(h)>0$ en $(x^*_-, x^*_+)$ y $sgn(h)<0$ en $(x^*_+,\infty).$
De esta forma es un máximo. Pero ¿es global? recordemos que necesitemos esa cota para todo $x\geq$.\\

Por el método llevado a cabo, $x^*_+$ es el único máximo local de $h$, que es continua, $h(0) = 0,\ \lim_{x\to \infty} = \lambda^2$ y además $h$ es creciente en $[0,x^*_+]$ y decreciente para $x>x^*_+$, por lo tanto $K = h(x^*_+)>0$ es el máximo en $[0,\infty)$ y por ende, sirve como nuestra cota.\\

Hemos concluido que se cumple la condición de crecimiento.



\subsubsection*{Alternativas para existencia y unicidad}
ver \cite{YamadaWatanabe1971}

\subsection*{Solución}
Podemos usar Lamperti, en particular usando $Y_t = \frac{2}{\sqrt{\lambda}}\sqrt{X_t}$

\begin{align*}
	dY_t &= \left(-\frac{\lambda}{\sqrt{\lambda}}\sqrt{X_t} + \frac{1}{2}\lambda X_t \cdot \left(-\frac{1}{2\sqrt{\lambda}}X_t^{-3/2}\right)\right)dt + dW_t \\
	&= \left(-\sqrt{\lambda}\sqrt{X_t} - \frac{\sqrt{\lambda}}{4\sqrt{X_t}}\right)dt + dW_t \\
	&= \left(-\frac{\lambda}{2}Y_t - \frac{1}{2Y_t}\right)dt + dW_t
\end{align*}

Ahora hay que intentar resolver esta ecuación con ruido aditivo.

[Tal vez \cite{MasoliverPerello2012} pueda ser útil.]

	

\printbibliography	

\end{document}
